% ============================================================
% SECCIÓN 6: SDK Científico de Python — Demeter SDK
% ============================================================
\chapter{SDK Científico de Python — \texttt{demeter\_sdk}}
\label{sec:sdk}

% ─────────────────────────────────────────────────────────────
\section{Propósito y Filosofía}

El \textbf{Demeter SDK} es la capa de acceso programático al sistema para investigadores y
analistas. Nacido de la necesidad de combinar los datos de telemetría del invernadero con
algoritmos agronómicos reconocidos, el SDK resuelve el ciclo completo de análisis en
\textit{una sola llamada}:

\begin{itemize}[leftmargin=*, label=--]
    \item Conectarse a la API del servidor con autenticación por clave.
    \item Descargar la telemetría de un experimento con caché automática.
    \item Limpiar y normalizar los datos (ETL).
    \item Calcular indicadores agronómicos derivados (VPD, GDD, ET$_0$, …).
    \item Exportar los resultados a CSV o Excel estructurado.
    \item Generar figuras de publicación científica con un único comando.
\end{itemize}

\begin{designnote}
\textbf{Filosofía de uso}: el investigador no necesita conocer la arquitectura interna del
sistema, ni Redis, ni WebSocket, ni la API directamente. Basta con instanciar
\texttt{DemeterClient}, pasar el ID del experimento y el SDK hace el resto.
El objetivo es que un científico con conocimientos básicos de Python pueda obtener
un informe completo en menos de diez líneas de código.
\end{designnote}

% ─────────────────────────────────────────────────────────────
\section{Arquitectura Modular}

El SDK está organizado en seis módulos con responsabilidades estrictamente separadas.
La clase \texttt{DemeterClient} actúa como orquestador público que los une:

\begin{table}[H]
    \centering
    \caption{Módulos del Demeter SDK}
    \begin{tabularx}{\textwidth}{llX}
        \toprule
        \textbf{Módulo}           & \textbf{Rol}              & \textbf{Clase / Función principal} \\
        \midrule
        \texttt{client.py}        & Orquestador               & \texttt{DemeterClient} \\
        \texttt{fetcher.py}       & Obtención de datos        & \texttt{fetch\_experiment()}, \texttt{fetch\_plant()} \\
        \texttt{transform.py}     & ETL (limpieza)            & \texttt{pipeline()}, \texttt{resample()} \\
        \texttt{science.py}       & Cálculos agronómicos      & \texttt{enrich()}, funciones individuales \\
        \texttt{viz.py}           & Visualización científica  & \texttt{plot\_timeseries()}, \texttt{plot\_vpd()}, … \\
        \texttt{export.py}        & Exportación de resultados & \texttt{to\_csv()}, \texttt{to\_excel()} \\
        \bottomrule
    \end{tabularx}
    \label{tab:sdk_modules}
\end{table}

% ─────────────────────────────────────────────────────────────
\section{Instalación y Configuración}

El SDK se instala como un paquete Python local desde la carpeta del proyecto:

\begin{lstlisting}[style=trama, caption={Instalación del SDK en modo editable}]
pip install -e ./SDK/
\end{lstlisting}

La autenticación se realiza mediante una \textbf{clave API} (\texttt{X-API-Key}) generada
por el sistema. Puede pasarse directamente al constructor o mediante la variable de entorno
\texttt{DEMETER\_API\_KEY}:

\begin{lstlisting}[style=trama, caption={Inicialización del cliente}]
from demeter_sdk import DemeterClient

# Opción 1: API key explícita
client = DemeterClient(
    api_key   = "dk_live_XXXXXXXXXXXX",
    base_url  = "https://demeter.tuservidor.com",
    cache_ttl = 3600,    # Caché de 1 hora
)

# Opción 2: variable de entorno DEMETER_API_KEY
client = DemeterClient(base_url="https://demeter.tuservidor.com")
\end{lstlisting}

% ─────────────────────────────────────────────────────────────
\section{Obtención de Datos: Módulo Fetcher}

El fetcher implementa una caché transparente en formato \textbf{Apache Parquet} para
evitar descargar los mismos datos repetidamente durante análisis iterativos. El comportamiento es:

\begin{itemize}[leftmargin=*, label=--]
    \item Si existe un fichero Parquet en \texttt{\~{}/.demeter\_cache/} con antigüedad
          inferior al TTL configurado (defecto: 1~h), se devuelve sin llamar a la API.
    \item Si el fichero es más antiguo o no existe, se descarga desde la API y se guarda.
    \item \texttt{force\_refresh=True} omite la caché siempre.
\end{itemize}

\begin{lstlisting}[style=trama, caption={Descarga de datos con caché automática}]
# Datos crudos de un experimento (últimos 30 días)
raw = client.get_raw_data(experimento_id=1, dias=30)

# Datos crudos de una planta individual (sirve el endpoint Redis)
raw_planta = client.get_plant_data(plant_id=7)

# Forzar recarga desde la API
raw_fresh = client.get_raw_data(experimento_id=1, force_refresh=True)
\end{lstlisting}

% ─────────────────────────────────────────────────────────────
\section{Pipeline ETL: Módulo Transform}

Antes de calcular cualquier indicador, los datos crudos pasan por un pipeline de limpieza
en cuatro etapas:

\begin{table}[H]
    \centering
    \caption{Etapas del pipeline ETL}
    \begin{tabularx}{\textwidth}{clX}
        \toprule
        \textbf{Paso} & \textbf{Función}   & \textbf{Qué hace} \\
        \midrule
        1 & \texttt{to\_dataframe()} & Convierte los JSON crudos a un DataFrame tipado. Parsea timestamps ISO-8601, coerce numérico de temperatura/humedad, ordena cronológicamente. \\
        2 & \texttt{clean()}         & Elimina filas duplicadas exactas y mediciones con ambos sensores a NaN (irrecuperables). \\
        3 & \texttt{fill\_gaps()}    & Interpolación lineal de huecos puntuales de sensor (desconexiones breves). Soporta también método \texttt{'cubic'} o \texttt{'time'}. \\
        4 & \texttt{resample()}      & Opcional. Agrega la serie temporal a menor frecuencia (e.g. \texttt{'1h'}, \texttt{'1D'}) por nodo. \\
        \bottomrule
    \end{tabularx}
\end{table}

\begin{lstlisting}[style=trama, caption={Uso del pipeline ETL}]
# Pipeline completo con resampleado horario
df = client.get_clean_data(
    experimento_id=1,
    dias=30,
    resample_rule="1h",   # Opcional
)
# Resultado: DataFrame con columnas [timestamp, temperature, humidity, node_id]
\end{lstlisting}

% ─────────────────────────────────────────────────────────────
\section{Cálculos Agronómicos: Módulo Science}

El módulo \texttt{science} implementa \textbf{todas las fórmulas] agronómicas} de forma
vectorizada sobre NumPy/Pandas, sin bucles Python. Todas las funciones admiten arrays o
\texttt{pd.Series} y devuelven \texttt{np.ndarray}.

\subsection{Termodinámica del Aire Húmedo}

\subsubsection{Presión de Vapor de Saturación ($e_s$)}

Calculada mediante la fórmula de \textbf{Magnus-Tetens}, el estándar de la meteorología
agrícola para el rango de temperaturas habitual en invernadero ($-10$ a $50$\,°C):

\[
e_s = 0{,}6108 \cdot e^{\displaystyle\frac{17{,}27 \cdot T}{T + 237{,}3}}
\]

\begin{description}
        \item [In:] $T$ — Temperatura de bulbo seco (°C).
    \item [Out:] $e_s$ — Presión de vapor de saturación (kPa).
\end{description}

\subsubsection{Déficit de Presión de Vapor (VPD)}

El VPD es el indicador más útil para el manejo del invernadero. Mide la diferencia entre
la cantidad de agua que el aire \textit{podría} contener ($e_s$) y la que \textit{realmente}
contiene ($e_a = e_s \cdot HR/100$):

\[
VPD = \max\left(0,\; e_s - e_a\right)
\]

\begin{description}
        \item [In:] $e_s$ — Presión de vapor de saturación (kPa).
    \item [In:] $e_a$ — Presión de vapor actual (kPa). Se calcula como $e_s \cdot HR / 100$.
    \item [Out:] $VPD$ — Déficit de presión de vapor (kPa).
\end{description}

La tabla siguiente define las zonas de riesgo para cultivos de invernadero típicos:

\begin{table}[H]
    \centering
    \caption{Zonas de riesgo del VPD para cultivos de invernadero}
    \begin{tabular}{clp{6cm}}
        \toprule
        \textbf{VPD (kPa)} & \textbf{Zona}          & \textbf{Implicación agronómica} \\
        \midrule
        $< 0{,}4$           & Riesgo fúngico         & Exceso de humedad; condensación en hoja probable. \\
        $0{,}4 - 1{,}0$     & Óptimo                 & Transpiración activa, absorción de nutrientes eficiente. \\
        $1{,}0 - 1{,}5$     & Estrés moderado        & Inicio de cierre estomático; reducir temperatura. \\
        $> 1{,}5$           & Estrés severo          & Cierre estomático; parar riego por goteo. \\
        \bottomrule
    \end{tabular}
    \label{tab:vpd_zones}
\end{table}

\subsubsection{Punto de Rocío ($T_d$)}

Temperatura a la que el aire se satura si se enfría isohúmamente. Cuando la temperatura
de una superficie (e.g. hoja) cae por debajo de $T_d$, hay condensación. Se calcula
invirtiendo la fórmula de Magnus:

\[
T_d = \frac{237{,}3 \cdot \ln(e_a / 0{,}6108)}{17{,}27 - \ln(e_a / 0{,}6108)}
\]

\begin{description}
        \item [In:] $e_a$ — Presión de vapor actual (kPa).
    \item [Out:] $T_d$ — Temperatura de punto de rocío (°C).
\end{description}

\subsubsection{Temperatura de Bulbo Húmedo ($T_w$) — Stull 2011}

Temperatura límite del enfriamiento por evaporación. Usada en sistemas de
refrigeración evaporativa de invernadero:

{\small
\[
\begin{aligned}
T_w = &\; T \cdot \arctan\!\left(0{,}151977 \cdot \sqrt{HR + 8{,}313659}\right) \\
      &\; + \arctan(T+HR) - \arctan(HR-1{,}676331) \\
      &\; + 0{,}00391838 \cdot HR^{1{,}5} \cdot \arctan(0{,}023101 \cdot HR) \\
      &\; - 4{,}686035
\end{aligned}
\]
}

\begin{description}
        \item [In:] $T$ — Temperatura de bulbo seco (°C).
    \item [In:] $HR$ — Humedad relativa (\%).
    \item [Out:] $T_w$ — Temperatura de bulbo húmedo (°C).
\end{description}

\subsubsection{Humedad Absoluta}

Masa de vapor de agua por metro cúbico de aire, independiente de la presión:
\[
AH = \frac{2165 \cdot e_a}{T + 273{,}15}
\]

\begin{description}
        \item [In:] $e_a$ — Presión de vapor actual (kPa).
    \item [In:] $T$ — Temperatura (°C).
    \item [Out:] $AH$ — Humedad absoluta (g/m$^3$).
\end{description}

\subsubsection{Entalpía Específica del Aire Húmedo — ASHRAE}

Energía total (térmica + latente) del aire húmedo por kilogramo de aire seco.
Clave para el dimensionado de sistemas de climatización y refrigeración:

\[
r = \frac{0{,}622 \cdot e_a}{P - e_a}
\]
\[
h = 1{,}006 \cdot T + r \cdot (2501 + 1{,}86 \cdot T)
\]

\begin{description}
        \item [In:] $T$ — Temperatura de bulbo seco (°C).
    \item [In:] $e_a$ — Presión de vapor actual (kPa).
    \item [In:] $P$ — Presión atmosférica (kPa, defecto 101,325).
    \item [Out:] $h$ — Entalpía específica del aire húmedo (kJ/kg aire seco).
    \item [Out:] $r$ — Relación de mezcla o humedad específica (kg agua/kg aire seco).
\end{description}

\subsubsection{Índice de Calor (Rothfusz, NWS)}

Temperatura \textit{percibida} por el ser humano (y por el cultivo) combinando temperatura
y humedad. Solo significativo para $T \geq 27$\,°C:

\begin{align*}
HI = &\;-8{,}785 + 1{,}611\,T + 2{,}339\,HR - 0{,}1461\,T\cdot HR \\
     &\; -0{,}01231\,T^2 - 0{,}01642\,HR^2 + 0{,}002212\,T^2\cdot HR \\
     &\; +0{,}000726\,T\cdot HR^2 - 0{,}000004\,T^2\cdot HR^2
\end{align*}

\begin{description}
        \item [In:] $T$ — Temperatura ambiental (°C).
    \item [In:] $HR$ — Humedad relativa (\%).
    \item [Out:] $HI$ — Índice de calor percibido (°C).
\end{description}

\subsubsection{Riesgo de Rocío en Hoja (LSD)}

Índice que evalúa si la superficie de la hoja (modelada a $T_{leaf} = T_{air} - 2$\,°C)
está por debajo de la temperatura de punto de rocío. Un valor negativo indica
condensación activa y riesgo de enfermedades fúngicas:

\[
LSD = e_s(T_{leaf}) - e_a(T_{air}, HR)
\]

\begin{description}
        \item [In:] $T_{leaf}$ — Temperatura de la hoja (modelada como $T_{air} - 2$, °C).
    \item [In:] $e_s$ — Presión de vapor de saturación a $T_{leaf}$ (kPa).
    \item [In:] $e_a$ — Presión de vapor actual del aire (kPa).
    \item [Out:] $LSD$ — Riesgo de rocío en hoja (\textit{Leaf Surface Dew}, kPa).
\end{description}

% ─────────────────────────────────────────────────────────────
\subsection{Fenología y Tiempo}

\subsubsection{Grados-Día de Crecimiento (GDD)}

El GDD es la unidad de acumulación de calor vegetal; establece el ritmo de desarrollo
fenológico del cultivo. Se calcula como:

\[
GDD = \max\!\left(0,\; \frac{T_{max} + T_{min}}{2} - T_{base}\right)
\]

\begin{description}
        \item [In:] $T_{max}$ — Temperatura máxima diaria (°C).
    \item [In:] $T_{min}$ — Temperatura mínima diaria (°C).
    \item [In:] $T_{base}$ — Temperatura umbral de crecimiento (°C, defecto 10).
    \item [Out:] $GDD$ — Grados-Día de Crecimiento (°C·día).
\end{description}

\subsubsection{Horas de Frío (\textit{Chill Hours})}

Recuento de horas en las que la temperatura se encuentra en el rango $(0, 7{,}2]$\,°C.
Indicador clave para predecir la ruptura de dormancia en frutales de hoja caduca.

\subsubsection{Evapotranspiración de Referencia — Hargreaves-Samani ($ET_0$)}

Estima la demanda hídrica de referencia usando solo datos de temperatura, sin
necesidad de piranómetros ni sensores de viento:

\[
ET_0 = 0{,}0023 \cdot R_a \cdot (T_{mean} + 17{,}8) \cdot \sqrt{T_{max} - T_{min}}
\]

\begin{description}
        \item [In:] $R_a$ — Radiación extraterrestre (equivalente mm/día, defecto 15).
    \item [In:] $T_{mean}$ — Temperatura media diaria (°C).
    \item [In:] $T_{max}$ — Temperatura máxima diaria (°C).
    \item [In:] $T_{min}$ — Temperatura mínima diaria (°C).
    \item [Out:] $ET_0$ — Evapotranspiración de referencia (mm/día).
\end{description}

\subsubsection{Días Tras la Siembra (DAP)}

Número de días enteros transcurridos desde la fecha de siembra:

\[
DAP = \left\lfloor t_{sample} - t_{siembra} \right\rfloor
\]

\begin{description}
        \item [In:] $t_{sample}$ — Fecha de la muestra (timestamp).
    \item [In:] $t_{siembra}$ — Fecha de siembra (timestamp).
    \item [Out:] $DAP$ — Días tras la siembra (días).
\end{description}

Permite contextualizar cada medición en términos del ciclo de vida del cultivo.

% ─────────────────────────────────────────────────────────────
\subsection{Estadística y Detección de Anomalías}

\subsubsection{Z-Score}

Normalización lineal de una serie temporal. Detecta cuántas desviaciones estándar
se aleja cada muestra de la media del periodo:

\[
z = \frac{x - \mu}{\sigma}
\]

\begin{description}
        \item [In:] $x$ — Valor de la muestra.
    \item [In:] $\mu$ — Media aritmética de la serie.
    \item [In:] $\sigma$ — Desviación estándar de la serie.
    \item [Out:] $z$ — Puntuación estándar (adim.).
\end{description}

El SDK aplica Z-Score de forma independiente a temperatura y humedad.
Se considera outlier cualquier muestra con $|z| > 3$ (umbral configurable).

\subsubsection{Índice de Riesgo Wallin (1962) — Tizón Tardío}

Modelo clásico de epidemiología vegetal para estimación del riesgo de \textit{Phytophthora
infestans} (mildiu/tizón tardío). Evalúa un periodo horario y devuelve un nivel de riesgo:

\begin{table}[H]
    \centering
    \caption{Lógica del índice de riesgo Wallin}
    \begin{tabular}{lll}
        \toprule
        \textbf{Condición} & \textbf{Nivel} & \textbf{Acción recomendada} \\
        \midrule
        Horas con HR $\geq 90$\% $< 10$            & 0 — Sin riesgo     & Sin acción. \\
        Horas $\geq 10$ y $10$°C $\leq T \leq 27$°C & 1 — Riesgo moderado & Revisión preventiva. \\
        Horas $\geq 20$                              & 2 — Riesgo alto    & Tratamiento fungicida. \\
        \bottomrule
    \end{tabular}
\end{table}

\subsubsection{Media Móvil 24h}

Suavizado de la serie temporal por grupos de nodo, usando una ventana de
24 muestras (equivalente a 24~h cuando los datos son horarios). Útil para
eliminar el ruido de alta frecuencia y visualizar tendencias.

% ─────────────────────────────────────────────────────────────
\section{El Helper \texttt{enrich()}: Enriquecimiento en Una Llamada}

La función \texttt{science.enrich()} aplica en un solo paso todos los cálculos
sobre el DataFrame limpio y devuelve un nuevo DataFrame con las columnas derivadas:

\begin{table}[H]
    \centering
    \caption{Columnas añadidas por \texttt{enrich()}}
    \begin{tabular}{lll}
        \toprule
        \textbf{Columna}         & \textbf{Unidad} & \textbf{Fórmula / Método} \\
        \midrule
        \texttt{vpd\_kpa}        & kPa             & Magnus-Tetens \\
        \texttt{dew\_point\_c}   & °C              & Inversión Magnus \\
        \texttt{wet\_bulb\_c}    & °C              & Stull 2011 \\
        \texttt{abs\_humidity\_g\_m3} & g/m$^3$   & Ley del gas ideal \\
        \texttt{enthalpy\_kj\_kg}& kJ/kg           & Definición ASHRAE \\
        \texttt{heat\_index\_c}  & °C              & Rothfusz (NWS) \\
        \texttt{lsd\_kpa}        & kPa             & Modelado hoja -2°C \\
        \texttt{z\_score\_temp}  & adim.           & $(T - \mu)/\sigma$ \\
        \texttt{z\_score\_hum}   & adim.           & $(HR - \mu)/\sigma$ \\
        \texttt{temp\_ma\_24h}   & °C              & Media móvil 24 muestras \\
        \texttt{hum\_ma\_24h}    & \%              & Media móvil 24 muestras \\
        \texttt{is\_outlier}     & bool            & $|z_{temp}| > 3{,}0$ \\
        \texttt{dap\_days}       & días            & Solo si \texttt{fecha\_siembra} provisto \\
        \bottomrule
    \end{tabular}
    \label{tab:enrich_columns}
\end{table}

\begin{lstlisting}[style=trama, caption={Pipeline completo con enriquecimiento}]
# Datos limpios + todos los indicadores agronómicos
df = client.get_enriched_data(
    experimento_id = 1,
    dias           = 30,
    fecha_siembra  = "2025-09-01",   # Activa columna dap_days
    resample_rule  = "1h",           # Opcional: resampleado horario
)

# Consultar el VPD medio de cada nodo
print(df.groupby("node_id")["vpd_kpa"].mean())

# Ver anomalías de temperatura
outliers = df[df["is_outlier"]]
print(f"{len(outliers)} lecturas anómalas detectadas")
\end{lstlisting}

% ─────────────────────────────────────────────────────────────
\section{Visualización Científica: Módulo Viz}

Todos los gráficos están implementados con Matplotlib y Seaborn con estilo
consistente y devuelven el objeto \texttt{Figure} para que el investigador pueda
añadir anotaciones, cambiar el título o guardar con el DPI deseado.

\begin{table}[H]
    \centering
    \caption{Funciones de visualización disponibles}
    \begin{tabular}{llp{7cm}}
        \toprule
        \textbf{Función}          & \textbf{Tipo}          & \textbf{Descripción} \\
        \midrule
        \texttt{plot\_timeseries()} & Serie temporal         & Cualquier métrica frente al tiempo, coloreada por nodo, con media móvil superpuesta. \\
        \texttt{plot\_vpd()}        & VPD con zonas          & VPD en el tiempo con bandas de color que marcan las 4 zonas de riesgo agronómico. \\
        \texttt{plot\_boxplot()}    & Diagrama de cajas      & Distribución de cualquier métrica por nodo; ideal para comparar variabilidad entre plantas. \\
        \texttt{plot\_heatmap()}    & Mapa de calor 2D       & Intensidad de la métrica en el tiempo (eje X) por nodo (eje Y); resampleado diario por defecto. \\
        \bottomrule
    \end{tabular}
\end{table}

\begin{lstlisting}[style=trama, caption={Generación de gráficos}]
# Serie temporal del VPD con zonas de riesgo
fig = client.plot.plot_vpd(df)
fig.savefig("vpd_zonas.png", dpi=300, bbox_inches="tight")

# Boxplot de temperatura por nodo
fig2 = client.plot.plot_boxplot(df, metric="temperature")
fig2.savefig("boxplot_temperatura.png", dpi=300)

# Heatmap de humedad relativa (dia x nodo)
fig3 = client.plot.plot_heatmap(df, metric="humidity", resample_rule="1D")
fig3.savefig("heatmap_humidity.png", dpi=300)
\end{lstlisting}

% ─────────────────────────────────────────────────────────────
\section{Exportación Científica: Módulo Export}

\subsection{CSV}

El fichero CSV se escribe en codificación \textbf{UTF-8 con BOM} (\texttt{utf-8-sig})
para asegurar la apertura correcta en Excel sin selección manual de codificación:

\begin{lstlisting}[style=trama, caption={Exportación a CSV}]
client.export.to_csv(df_enriched, "resultados/telemetria.csv")
\end{lstlisting}

\subsection{Excel Multi-hoja}

El libro Excel generado tiene tres hojas estructuradas para la comunicación científica:

\begin{table}[H]
    \centering
    \caption{Estructura del libro Excel generado}
    \begin{tabularx}{\textwidth}{lXX}
        \toprule
        \textbf{Hoja}              & \textbf{Contenido} & \textbf{Utilidad} \\
        \midrule
        \texttt{Resumen}           & Media, mediana, desv. típica, mín., máx., P5, P95 de todas las métricas. & Vista rápida para el investigador principal y presentaciones. \\
        \texttt{Datos\_Enriquecidos} & Todos los registros con todas las columnas derivadas (13+ columnas). & Dataset completo para análisis estadístico externo (R, SPSS). \\
        \texttt{Datos\_Crudos}     & Solo temperatura y humedad originales (sin derivados). & Reproducibilidad: los datos tal y como llegaron del sensor. \\
        \bottomrule
    \end{tabularx}
\end{table}

\begin{lstlisting}[style=trama, caption={Exportación a Excel estructurado}]
client.export.to_excel(
    df_enriched,
    "resultados/informe.xlsx",
    df_raw             = df_raw,
    experiment_name    = "Ensayo Estrés Salino 2025",
)
\end{lstlisting}

% ─────────────────────────────────────────────────────────────
\section{Suite Completa — Un Solo Comando}

Para el flujo de análisis estándar de un experimento, el método \texttt{run\_full\_suite()}
ejecuta automáticamente: descarga, ETL, enriquecimiento, exportación CSV + Excel y los
cuatro tipos de gráfico, guardándolo todo en una estructura de carpetas organizada:

\begin{lstlisting}[style=trama, caption={Análisis completo en una sola llamada}]
ruta = client.run_full_suite(
    experimento_id = 3,
    dias           = 60,
    output_dir     = "./resultados",
    fecha_siembra  = "2025-10-01",
    resample_rule  = "1h",
)
# Resultado:
# resultados/Experimento_3/
# ├── data/
# │   ├── enriched_telemetry.csv
# │   └── report.xlsx
# └── plots/
#     ├── temperature_timeseries.png
#     ├── vpd_zones.png
#     ├── temperature_boxplot.png
#     └── heatmap_temperature.png
print(f"Análisis guardado en: {ruta}")
\end{lstlisting}

% ─────────────────────────────────────────────────────────────
\section{Ejemplo de Sesión Completa con Jupyter}

El siguiente extracto reproduciría una sesión típica de investigación en un Jupyter Notebook,
desde la conexión hasta la interpretación de resultados:

\begin{lstlisting}[style=trama, caption={Sesión completa en Jupyter Notebook}]
from demeter_sdk import DemeterClient

# 1. Conectar
client = DemeterClient(api_key="dk_live_XXX",
                       base_url="https://demeter.servidor.com")
assert client.ping(), "Servidor no disponible"

# 2. Obtener datos limpios y enriquecidos
df = client.get_enriched_data(
    experimento_id=2,
    dias=30,
    fecha_siembra="2025-09-15",
    resample_rule="30min",
)

# 3. Consultas rápidas
vpd_alto = df[df["vpd_kpa"] > 1.5]
print(f"Horas con estrés severo: {len(vpd_alto)}")

wallin = client.plot.plot_vpd(df)
wallin.savefig("vpd.png", dpi=300)

# 4. Calcular GDD acumulado del experimento
from demeter_sdk.science import gdd
df_diario = df.resample("1D", on="timestamp").agg(
    T_max=("temperature", "max"),
    T_min=("temperature", "min")
).reset_index()
df_diario["gdd"] = gdd(df_diario["T_max"], df_diario["T_min"])
df_diario["gdd_acum"] = df_diario["gdd"].cumsum()
print(df_diario[["timestamp", "gdd_acum"]].tail())

# 5. Exportar
client.export.to_excel(df, "informe_exp2.xlsx", experiment_name="Exp 2")
\end{lstlisting}

% ─────────────────────────────────────────────────────────────
\section{Referencia Rápida de Parámetros}

\begin{table}[H]
    \centering
    \caption{Parámetros principales del \texttt{DemeterClient}}
    \begin{tabularx}{\textwidth}{llX}
        \toprule
        \textbf{Parámetro}  & \textbf{Tipo}  & \textbf{Descripción} \\
        \midrule
        \texttt{api\_key}       & str    & Clave de acceso a la API (o via \texttt{DEMETER\_API\_KEY}). \\
        \texttt{base\_url}      & str    & URL base del servidor (defecto: \texttt{localhost:8000}). \\
        \texttt{cache\_dir}     & str    & Ruta de la caché Parquet (defecto: \texttt{\~{}/.demeter\_cache}). \\
        \texttt{cache\_ttl}     & int    & Tiempo de vida en segundos de los ficheros Parquet (defecto: 3600). \\
        \texttt{timeout}        & int    & Timeout por petición HTTP en segundos (defecto: 30). \\
        \texttt{max\_retries}   & int    & Reintentos automáticos ante errores 5xx (defecto: 3). \\
        \bottomrule
    \end{tabularx}
\end{table}
